\BeleskeID{02-s049b}
% element: 02-s049b:d01

Pri zameni troulaznog NILI kola pratimo međusignale:
\[
u=\overline{X+Y},\qquad v=\overline{u+u}=X+Y,\qquad
w=\overline{v+Z}=\overline{X+Y+Z}.
\]
NILI kolo sa spojenim ulazima daje negaciju signala. Prvo kolo i invertor ponovo daju pravi zbir \(X+Y\), a završno kolo negira zbir sa \(Z\). Zato nije ispravno neposredno kaskadirati samo dva NILI kola bez međuinverzije. Dodatni nivoi menjaju vreme odziva iako je stabilna funkcija ista.

U tri četvorostruka kućišta kupujemo ukupno 12 kola, od kojih koristimo deset, pa dva ostaju slobodna. Razlikujemo broj upotrebljenih kola od broja kupljenih kućišta.

Zadatak poređenja jednostepenih složenih CMOS realizacija polazi od dostupnih komplementnih ulaza. Treba sastaviti odgovarajuće tranzistorske mreže, utvrditi serijske i paralelne putanje i uporediti potrebne dimenzije za iste zahteve pogona. Broj literala ili kućišta sam ne daje površinu takvog kola; zaključak treba zasnovati na konkretnoj CMOS realizaciji.
